%% ======================================================================
%%  Springer Nature journal paper, TWO-COLUMN (sn-jnl class, template v3.1, Dec 2024)
%%
%%  Build:  pdflatex Springer_paper_twocolumn -> bibtex Springer_paper_twocolumn -> pdflatex x2
%%  Overleaf: upload this folder (Springer_paper_twocolumn.tex, sn-jnl.cls,
%%            sn-mathphys-num.bst, references.bib, figures/),
%%            compiler = pdfLaTeX.
%%
%%  Content is the same study as ../IEEE/IEEE_paper.tex. Every number is
%%  produced by the scripts in ../IEEE/scripts/ from the trained model
%%  models/taka_authenticator.pth and the public JaalTaka dataset.
%% ======================================================================

%% Two-column Springer Nature layout (option iicol). Single-column version:
%% Springer_paper.tex. For line numbers, add the "lineno" option.
\documentclass[pdflatex,iicol,sn-mathphys-num]{sn-jnl}

\usepackage{graphicx}%
\usepackage{multirow}%
\usepackage{amsmath,amssymb,amsfonts}%
\usepackage{xcolor}%
\usepackage{textcomp}%
\usepackage{booktabs}%

\graphicspath{{figures/}}
\hyphenation{bank-note bank-notes counter-feit counter-feits Efficient-Net}

\raggedbottom

\begin{document}

\title[Reliable Smartphone-Based Detection of Counterfeit Bangladeshi Banknotes]{Reliable
Smartphone-Based Detection of Counterfeit Bangladeshi Banknotes: Multi-View Deep Learning
with Out-of-Distribution Rejection and a Benchmark Audit}

\author*[1]{\fnm{Debobrato} \sur{Biswas}}\email{biswas2305343003@diu.edu.bd}

\author[1]{\fnm{Imran} \sur{Mahmud}}

\affil*[1]{\orgdiv{Department of Software Engineering},
\orgname{Daffodil International University},
\orgaddress{\city{Dhaka}, \postcode{1207}, \country{Bangladesh}}}

%% Structured abstract (Springer Nature style: no citations or equations)
\abstract{\textbf{Background:} Counterfeit high-value banknotes harm retailers and
individuals who cannot afford verification hardware.

\textbf{Methods:} We present a smartphone-based authenticator for Bangladeshi Tk~500 and
Tk~1000 notes, evaluated on JaalTaka (1{,}390 notes: 802 genuine, 588 seized
counterfeits; six photographs each). It fine-tunes EfficientNet-B0, fuses up to six
photographs per note, returns a genuine, uncertain, or counterfeit verdict, rejects
unsupported inputs, and explains decisions with Grad-CAM. We audited the benchmark for
duplicates.

\textbf{Results:} On 208 held-out test notes, the system reaches 98.08\% accuracy (95\%
CI 95.16--99.25\%) and an ROC-AUC of 0.994 at 37~ms per photograph. However, the 1{,}390
notes form only 846 distinct photograph sets, 116 test notes have a copy in training, and
12 duplicate groups carry both labels, causing all four test errors. On the 92 test notes
without a copy, the system and two simple baselines are error-free.

\textbf{Conclusion:} The benchmark is saturated. We release the duplicate groups and
recommend group-aware, counterfeit-family hold-out evaluation.}

\keywords{Banknote authentication, Bangladeshi Taka, Counterfeit detection, Data leakage,
EfficientNet, Explainable AI, Grad-CAM, Out-of-distribution detection, Transfer learning}

\maketitle

%% Plain Language Summary and Word Count, as requested by manuscript screening.
\section*{Plain Language Summary}
Fake Tk~500 and Tk~1000 notes are hard to spot by eye. We built a tool that checks a
banknote from one to six phone photographs, answers ``genuine'', ``counterfeit'', or
``uncertain'', refuses pictures of other objects, and highlights the parts of the note
behind its answer. It was correct for 98\% of the test notes of a public photo
collection. That collection, however, contains many repeated photographs, some labelled
both genuine and fake, and these explained every mistake. The tool may therefore not
catch new kinds of fakes. We share the repeated photographs and suggest fairer tests.

\section*{Word Count}
2{,}842 words excluding references (main text: 1{,}944 words; abstract: 157 words);
9 figures; 4 tables.

%% ======================================================================
\section{Introduction}\label{sec:intro}

Counterfeit Tk~500 and Tk~1000 notes undermine trust in cash, the main means of payment
for small retailers in Bangladesh. Banks use ultraviolet, infrared, and magnetic sorters,
but shopkeepers and citizens rely on visual and tactile checks that are slow and
unreliable for untrained users. A smartphone camera, carried by almost everyone, makes
image-based authentication an attractive complement.

The JaalTaka dataset~\cite{piyas2025jaaltaka} enabled such work for Bangladeshi notes by
releasing smartphone photographs of genuine notes and of counterfeits seized by the Rapid
Action Battalion (RAB); its authors reported 97.4--98.9\% accuracy for five ImageNet
CNNs. High benchmark accuracy is necessary but not sufficient for a tool that people
trust with money. It must also (i)~handle unsupported inputs, (ii)~admit uncertainty,
(iii)~explain its decisions, and (iv)~be evaluated in a way that predicts behaviour on
unseen notes. The last point is easy to get wrong: networks exploit
shortcuts~\cite{geirhos2020shortcut}, and duplicates across splits inflate
accuracy~\cite{kapoor2023leakage}. Our contributions are:
\begin{enumerate}
\item an open, CPU-real-time pipeline~(Fig.~\ref{fig:pipeline}) combining
EfficientNet-B0~\cite{tan2019efficientnet}, fusion of up to six photographs, a three-way
decision with adjustable strictness, an out-of-distribution (OOD) gate, and
Grad-CAM~\cite{selvaraju2017gradcam};
\item a note-disjoint evaluation with confidence intervals, threshold and OOD analyses,
and frozen-feature baselines; and
\item a duplicate audit showing that label conflicts between copied notes cause every
test error and that the benchmark is saturated, plus an occlusion test revealing signal
outside the note's centre.
\end{enumerate}

\begin{figure*}[!t]
\centering
\includegraphics[width=\textwidth]{fig_pipeline}
\caption{Overview of the system. Each photograph is classified by a fine-tuned
EfficientNet-B0; photographs failing the OOD test are excluded, the remaining genuine
probabilities are averaged, and a three-way verdict is returned with Grad-CAM maps}
\label{fig:pipeline}
\end{figure*}

%% ======================================================================
\section{Related Work}\label{sec:related}

Early detectors used hand-crafted features: luminance histograms with multiple-kernel
SVMs~\cite{yeh2011}, scanned background printing texture~\cite{berenguel2016}, the
multispectral response of inks~\cite{baek2018}, and individual security features of
Bangladeshi notes~\cite{uddin2016}; a mobile application was also
proposed~\cite{dittimi2018}. These methods depend on controlled imaging and
denomination-specific tuning. Transfer learning from
ImageNet~\cite{deng2009imagenet,yosinski2014transferable} lets CNNs learn from few
images: Pham et~al. detected fake banknotes in smartphone images~\cite{pham2020,pham2022},
Bangladeshi denominations were recognised with neural networks and lightweight
CNNs~\cite{jahangir2007,linkon2020}, JaalTaka was benchmarked~\cite{piyas2025jaaltaka}
with DenseNet121~\cite{huang2017densenet}, ResNet50~\cite{he2016resnet},
VGG16~\cite{simonyan2015vgg}, MobileNetV2~\cite{sandler2018mobilenetv2}, and
InceptionV3~\cite{szegedy2016inception}, and transfer learning was also applied
in~\cite{tabassum2025}. None of these studies evaluates OOD inputs, uncertainty, or
duplicates.

Maximum softmax probability is a weak OOD signal~\cite{hendrycks2017baseline};
Mahalanobis distances on network features are stronger~\cite{lee2018mahalanobis}, and we
adapt them to a low-dimensional subspace. Miscalibrated
probabilities~\cite{guo2017calibration} motivate an explicit ``uncertain'' outcome, and
Grad-CAM~\cite{selvaraju2017gradcam} shows the regions supporting a class. Shortcut
learning~\cite{geirhos2020shortcut} and train--test leakage~\cite{kapoor2023leakage}
motivate our audit.

%% ======================================================================
\section{Materials and Methods}\label{sec:method}

\subsection{Dataset}\label{sec:data}

JaalTaka~\cite{piyas2025jaaltaka,jaaltaka_data} (CC~BY~4.0) contains 1{,}390 notes:
802 genuine (398 Tk~500, 404 Tk~1000) from three national banks and 588 counterfeit
(320 Tk~500, 268 Tk~1000) supplied by RAB. Each note was photographed six times, showing
different segments~(Fig.~\ref{fig:samples}), with one of four smartphones (Samsung Galaxy
A50 and M31, Redmi Note~8 and Note~9~Pro), giving 8{,}340 images, which we resized to a
longer side of 1{,}024 pixels. According to the dataset description, 649 original notes
(502 genuine, 147 counterfeit) were expanded to 1{,}390 by augmentation.

\begin{figure*}[!t]
\centering
\includegraphics[width=\textwidth]{fig_samples}
\caption{The six photographs of a genuine Tk~1000 note (top) and a counterfeit Tk~500
note (bottom) from the test split; colour, lighting, and background vary with the phone
and session}
\label{fig:samples}
\end{figure*}

\subsection{Evaluation protocol}\label{sec:protocol}

Because the photographs of a note are correlated, we split by note: within each class,
notes were shuffled (seed 42) and assigned 70/15/15\% to training (562 genuine, 412
counterfeit; 5{,}844 photographs), validation, and test (120 genuine, 88 counterfeit;
1{,}248 photographs each). Validation serves model selection and OOD calibration; the
test split is used once. Section~\ref{sec:audit} shows that this note-disjoint split is
\emph{not} sufficient, because distinct notes can share photographs. Genuine is the
positive class. We report accuracy with a 95\% Wilson
interval~\cite{wilson1927probable}, the counterfeit detection rate (specificity), the
genuine acceptance rate (recall), and ROC-AUC~\cite{fawcett2006roc}, per photograph
($n=1{,}248$) and per note ($n=208$), at threshold 0.5 unless stated otherwise.

\subsection{Classifier and training}\label{sec:training}

EfficientNet-B0~\cite{tan2019efficientnet}, chosen for its accuracy--cost trade-off on
CPUs and initialised with ImageNet weights~\cite{deng2009imagenet}, maps a
$288\times288$ input $x$ to a $1280\times9\times9$ feature map; global average pooling
yields the embedding $\mathbf{e}=f(x)\in\mathbb{R}^{1280}$. The head is dropout
($p=0.3$) and a linear layer:
\begin{equation}
p(x) = \operatorname{softmax}\!\left(\mathbf{W}\mathbf{e}+\mathbf{b}\right)_{\text{genuine}} .
\label{eq:prob}
\end{equation}
The model has 4.01~M parameters. Training uses random resized crops (scale 0.35--1.0,
aspect 0.75--1.33), rotations up to $\pm10^\circ$, colour jitter (brightness and contrast
$\pm0.4$, saturation $\pm0.3$, hue $\pm0.04$), and Gaussian blur ($p=0.2$). At
evaluation, the shorter side is resized to 331 pixels and the central $288\times288$
crop is kept. The stem and first three MBConv stages are frozen with their ImageNet
batch-norm statistics; the rest is fine-tuned with the class-weighted,
label-smoothed~\cite{szegedy2016inception} cross-entropy
\begin{align}
\mathcal{L} &= -\sum_{c\in\{0,1\}} w_c\, q_c \log \hat{y}_c, \nonumber\\
q_c &= (1-\varepsilon)\,\mathbb{1}[c=y] + \tfrac{\varepsilon}{2},
\label{eq:loss}
\end{align}
with $\varepsilon=0.05$ and $w_c=N/(2N_c)$. We use AdamW~\cite{loshchilov2019adamw}
(weight decay $10^{-4}$), a one-cycle schedule~\cite{smith2019superconvergence} (peak
$3\times10^{-4}$, 15\% warm-up), batch size 12, and six epochs, keeping the epoch with
the highest validation accuracy plus AUC; accuracy rose from 98.80\% after the first
epoch to 99.04\% (AUC 0.9976) after the sixth, which was selected. Training took about
55~min on an Intel Core i5-12400 CPU with PyTorch~\cite{paszke2019pytorch}.

\subsection{Multi-view fusion and three-way decision}\label{sec:decision}

Let $\mathcal{V}$ be the photographs of a note that pass the OOD gate. The note-level
probability is the mean, a late-fusion form of multi-view pooling~\cite{su2015mvcnn},
\begin{equation}
\bar{p} = \frac{1}{|\mathcal{V}|}\sum_{i\in\mathcal{V}} p(x_i),
\label{eq:fusion}
\end{equation}
and, for a strictness threshold $\tau\in[0.5,1)$ (default 0.7), the verdict is
\begin{equation}
\text{verdict} =
\begin{cases}
\text{not supported} & \text{if } \mathcal{V}=\emptyset,\\
\text{genuine} & \text{if } \bar{p}\ge\tau,\\
\text{counterfeit} & \text{if } \bar{p}<0.5,\\
\text{uncertain} & \text{otherwise.}
\end{cases}
\label{eq:decision}
\end{equation}
Averaging needs no extra training and accepts any number of views. The asymmetry
reflects costs: accepting a counterfeit causes a loss, whereas ``uncertain'' only asks
the user to check the note by hand or add photographs. Users can choose
$\tau\in\{0.6,0.7,0.85\}$ or any other value.

\subsection{Out-of-distribution gate}\label{sec:ood}

Following~\cite{lee2018mahalanobis}, let $\boldsymbol{\mu}$ be the mean training
embedding and $\mathbf{U}_k$ its first $k=64$ principal directions with variances
$\sigma_j^2$. With $\mathbf{z}=\mathbf{U}_k^{\top}(\mathbf{e}-\boldsymbol{\mu})$, we
compute the in-subspace Mahalanobis distance and the out-of-subspace residual
\begin{equation}
m(\mathbf{e}) = \sum_{j=1}^{k}\frac{z_j^2}{\sigma_j^2},\qquad
\rho(\mathbf{e}) = \left\lVert(\mathbf{e}-\boldsymbol{\mu})-\mathbf{U}_k\mathbf{z}\right\rVert_2^2 .
\label{eq:ood}
\end{equation}
Thresholds $\theta_m$ and $\theta_\rho$ are the 99.5th validation percentiles. A
photograph is accepted when $r(\mathbf{e})=\max\{m/\theta_m,\ \rho/\theta_\rho\}\le1$
and its Laplacian variance (a texture measure) is at least 30, which rejects blank and
blurred images. The residual catches embeddings that leave the banknote subspace, which
the in-subspace distance alone cannot detect.

\subsection{Explanation and deployment}\label{sec:deploy}

Grad-CAM~\cite{selvaraju2017gradcam} is computed on the last feature map
$A\in\mathbb{R}^{1280\times9\times9}$ for the predicted class $c$:
\begin{equation}
\alpha_k = \frac{1}{81}\sum_{u,v}\frac{\partial y^{c}}{\partial A^{k}_{uv}},\qquad
L^{c} = \operatorname{ReLU}\!\Big(\sum_{k}\alpha_k A^{k}\Big),
\label{eq:gradcam}
\end{equation}
where $y^c$ is the logit of class $c$; the map is normalised, upsampled, and overlaid on
the photograph. The model runs in 37~ms per photograph (median of 30 runs, six CPU
threads). The deployed Flask and Streamlit interfaces also average a full-frame crop;
all reported results use the central crop only.

%% ======================================================================
\section{Results}\label{sec:results}

\subsection{Main results and baselines}\label{sec:main}

Per note, the system rejects 86 of 88 counterfeits and accepts 118 of 120 genuine notes~(Table~\ref{tab:main},~Fig.~\ref{fig:confusion}), and both ROC curves in~Fig.~\ref{fig:roc} rise almost vertically. The per-note confusion matrix equals the
per-photograph matrix divided by six, and each view position alone reaches 98.08\%:
every error comes from all six photographs of one of four notes.

\begin{table*}[!t]
\caption{Test results of the proposed system (threshold 0.5)}\label{tab:main}
\begin{tabular*}{\textwidth}{@{\extracolsep\fill}lcc}
\toprule
Metric & Per photograph & Per note (6 views) \\
\midrule
Samples (genuine / counterfeit) & 1{,}248 (720 / 528) & 208 (120 / 88) \\
Accuracy (\%) & 98.08 & 98.08 \\
\quad 95\% Wilson CI (\%) & 97.15--98.70 & 95.16--99.25 \\
Counterfeit detection (\%) & 97.73 & 97.73 (86/88) \\
Genuine acceptance (\%) & 98.33 & 98.33 (118/120) \\
Balanced accuracy (\%) & 98.03 & 98.03 \\
F1, genuine class & 0.983 & 0.983 \\
ROC-AUC & 0.990 & 0.994 \\
\botrule
\end{tabular*}
\end{table*}

\begin{figure}[!t]
\centering
\includegraphics[width=\columnwidth]{fig_confusion}
\caption{Test confusion matrices per photograph (left) and per note (right)}
\label{fig:confusion}
\end{figure}

\begin{figure}[!t]
\centering
\includegraphics[width=\columnwidth]{fig_roc}
\caption{Test ROC curves for per-note fusion, single photographs, and the occlusion
diagnostics of Section~\ref{sec:shortcut}}
\label{fig:roc}
\end{figure}

We also compared a cosine 1-nearest-neighbour (1-NN) classifier and a
logistic-regression probe on \emph{frozen} ImageNet EfficientNet-B0 embeddings~(Table~\ref{tab:baselines}).
Fine-tuning adds 1.4--1.6 points of per-photograph accuracy, but at note level the probe
matches it and 1-NN errs on only six notes. The JaalTaka authors' results
(97.4--98.9\%~\cite{piyas2025jaaltaka}) used a different split and are not directly
comparable.

\begin{table*}[!t]
\caption{Comparison with frozen-feature baselines on the test split}\label{tab:baselines}
\begin{tabular*}{\textwidth}{@{\extracolsep\fill}lcccc}
\toprule
& \multicolumn{2}{@{}c@{}}{Per photograph} & \multicolumn{2}{@{}c@{}}{Per note} \\
\cmidrule{2-3}\cmidrule{4-5}
Method & Acc.\ (\%) & AUC & Acc.\ (\%) & AUC \\
\midrule
1-NN, frozen ImageNet features & 96.71 & 0.967 & 97.12 & 0.979 \\
Linear probe, frozen features  & 96.47 & 0.981 & 98.08 & 0.983 \\
Fine-tuned EfficientNet-B0 (ours) & \textbf{98.08} & \textbf{0.990} & \textbf{98.08} & \textbf{0.994} \\
\botrule
\end{tabular*}
\end{table*}

\subsection{Fusion and the strictness threshold}\label{sec:threshold}

Over all $\binom{6}{k}$ subsets of $k$ views, accuracy stays at 98.08\% while the mean
AUC rises monotonically from 0.9900 ($k=1$) through 0.9912, 0.9917, 0.9925, and 0.9930
to 0.9939 ($k=6$): fusion makes the ranking of notes more reliable. The fused scores are
strongly bimodal~(Fig.~\ref{fig:scores}): correctly classified genuine notes have
$\bar{p}>0.86$ and counterfeits $\bar{p}<0.10$, so outcomes are unchanged for
$\tau\le0.85$, and at $\tau=0.90$ one genuine note (0.8\%) becomes ``uncertain''. The two
accepted counterfeits ($\bar{p}\approx0.95$) would need $\tau>0.954$, which would send
most genuine notes to manual inspection.

\begin{figure*}[!t]
\centering
\includegraphics[width=\textwidth]{fig_scores}
\caption{\textbf{a} Fused genuine probability $\bar{p}$ of the 208 test notes.
\textbf{b} Share of test notes accepted, ``uncertain'', or wrongly accepted versus the
threshold $\tau$}
\label{fig:scores}
\end{figure*}

\subsection{Out-of-distribution gate}\label{sec:oodres}

The gate rejects 4 of 1{,}248 test photographs (0.32\%) versus 0.72\% on validation,
with a median test ratio of 0.26~(Fig.~\ref{fig:ood}); no test note was ``not
supported''. In a preliminary check, it rejected 100 of 105 photographs of other
denominations (Tk~2--200) and blank, noise, and unrelated images; a systematic OOD
benchmark remains future work.

\begin{figure}[!t]
\centering
\includegraphics[width=\columnwidth]{fig_ood}
\caption{OOD ratio $r(\mathbf{e})$ of validation and test photographs; the dashed line
is the rejection limit}
\label{fig:ood}
\end{figure}

\subsection{Visual explanations}\label{sec:gradcam}

In~Fig.~\ref{fig:gradcam}, Grad-CAM concentrates on the portrait and the denomination
and ornament panels for correct decisions, and also on the security band for the
counterfeit in row~b. The wrongly accepted note in row~c looks genuine to the model, as
expected because an identical copy appears in training labelled genuine
(Section~\ref{sec:audit}).

\begin{figure*}[!t]
\centering
\includegraphics[width=\textwidth]{fig_gradcam}
\caption{Grad-CAM on test notes (views~1 and~2). \textbf{a} Genuine, accepted.
\textbf{b} Counterfeit, rejected. \textbf{c} Counterfeit, wrongly accepted; its
photographs occur in training under the genuine label}
\label{fig:gradcam}
\end{figure*}

\subsection{Shortcut diagnostic}\label{sec:shortcut}

As shown in~Fig.~\ref{fig:shortcut}, we blacked out either the centre (keeping a 20\%
border) or the border of each test input. Hiding the border gives 88.94\% accuracy (AUC
0.957). Hiding the centre gives 75.48\%, as most genuine photographs are rejected
(acceptance 58.47\%), yet the AUC remains 0.959 and 98.67\% of counterfeits are
recognised. The border, which holds the note's margins but often also the table or hand
beneath it, thus suffices to rank most photographs. Because genuine and counterfeit
notes came from different sources (banks and RAB), capture conditions may contribute to
decisions~\cite{geirhos2020shortcut}.

\begin{figure*}[!t]
\centering
\includegraphics[width=\textwidth]{fig_shortcut}
\caption{Occlusion diagnostic. \textbf{a}--\textbf{c} Inputs with the centre or the
20\% border hidden. \textbf{d} Per-photograph test results under each condition}
\label{fig:shortcut}
\end{figure*}

\subsection{Benchmark audit: duplicates and label conflicts}\label{sec:audit}

The perfect agreement between a note's photographs, the strong frozen baselines, and the
documented augmentation led us to search for duplicates. For every pair of notes, we
computed the normalised cross-correlation (NCC) of $48\times48$ grey-scale thumbnails of
corresponding views, averaged over the six views and maximised over flips. The
best-match distribution is bimodal~(Fig.~\ref{fig:leakage}a): for validation and test
notes, no match lies between 0.963 and 0.997. We therefore grouped notes with
NCC~$>0.98$~(Table~\ref{tab:audit}).

\begin{figure*}[!t]
\centering
\includegraphics[width=\textwidth]{fig_leakage}
\caption{Duplicate audit. \textbf{a} Best NCC match of each test note to any training
note; dashed line: threshold 0.98. \textbf{b}--\textbf{d} A test note (top) and its
matching training note (bottom); green frame: genuine, red: counterfeit. \textbf{b}
Augmented copy, same label. \textbf{c}, \textbf{d} Identical photographs with opposite
labels}
\label{fig:leakage}
\end{figure*}

\begin{table}[!t]
\caption{Duplicate audit of JaalTaka (pixel NCC $>0.98$ over six views)}\label{tab:audit}
\footnotesize
\begin{tabular*}{\columnwidth}{@{\extracolsep\fill}lr}
\toprule
Quantity & Value \\
\midrule
Notes / distinct photograph sets & 1{,}390 / 846 \\
Groups of size 1 / 2 / 3 & 449 / 250 / 147 \\
Notes sharing photographs with another & 941 (67.7\%) \\
Counterfeit-only / genuine-only groups & 289 / 545 \\
Groups containing \emph{both} labels & 12 (31 notes) \\
\midrule
Test notes with a copy in train/val & 116/208 (55.8\%) \\
\quad counterfeit / genuine & 60/88 / 56/120 \\
Test notes in label-conflict groups & 6 \\
Test errors due to label conflicts & 4 of 4 \\
\botrule
\end{tabular*}
\end{table}

The augmented copies were thus stored as independent notes, and a note-level split
places copies of 55.8\% of test notes in training or validation. Twelve groups contain
genuine- and counterfeit-labelled copies of the \emph{same} photographs (e.g., genuine
notes 400--406 duplicate counterfeit notes 106--112 and 547--553, and counterfeit notes
400 and 403--406 duplicate genuine notes 802 and 1--4), so at least one label per group
is wrong; genuine-labelled notes 402 and 403 carry serial number 1843834, which also
appears on counterfeit-labelled notes. \emph{All four} test errors are notes from these
groups, each predicted with the label of its training copy. On the 92 \emph{clean} test
notes, whose group has no member in training or validation, and on the 202 notes outside
label-conflict groups, all three methods are error-free~(Table~\ref{tab:clean}). The
98.08\% accuracy is thus limited by label noise, not by the model, and the task is
saturated even for 1-NN on generic features.

\begin{table*}[!t]
\caption{Note-level accuracy (\%) on audited test subsets}\label{tab:clean}
\begin{tabular*}{\textwidth}{@{\extracolsep\fill}lccc}
\toprule
Test subset & All & Without label conflicts & Clean (no copy) \\
\midrule
Notes (genuine / counterfeit) & 208 (120/88) & 202 (117/85) & 92 (64/28) \\
\midrule
1-NN, frozen features & 97.12 & 100.0 & 100.0 \\
Linear probe, frozen & 98.08 & 100.0 & 100.0 \\
Ours (fine-tuned) & 98.08 & 100.0 & 100.0 \\
\quad 95\% Wilson CI & 95.16--99.25 & 98.13--100 & 95.99--100 \\
\botrule
\end{tabular*}
\end{table*}

%% ======================================================================
\section{Discussion}\label{sec:discussion}

Within JaalTaka, the system is error-free once label conflicts are removed (lower 95\%
bound 96.0\% on the 92 clean notes), its errors are fully explained, and it runs in real
time on a CPU. This does not show that it will detect \emph{new} counterfeits: the 588
counterfeits collapse into 289 photograph sets, those we inspected share a few serial
numbers, and they were collected separately from the genuine notes. A model, even 1-NN,
can therefore succeed by recognising known batches or capture sessions, as the occlusion
test suggests. The saturated benchmark also cannot rank methods, since differences
between published CNNs are smaller than its label noise.

We recommend (i)~removing or relabelling the 12 conflict groups; (ii)~splitting by
duplicate group, for which we release our list; (iii)~holding out counterfeit families
(e.g., by serial number or seizure batch) and capture devices, to measure generalisation
to unseen forgeries and phones; and (iv)~reporting confidence intervals and
frozen-feature baselines, which immediately reveal saturation.

\subsection{Limitations}\label{sec:limits}

The system covers only Tk~500 and Tk~1000. Visible-light photographs cannot capture
ultraviolet, magnetic, or tactile features, so the output is decision support and
suspicious notes should be confirmed at a bank. Scans and distant photographs are mostly
rejected, the OOD evaluation is preliminary, and probabilities are not calibrated on
independent data~\cite{guo2017calibration}. Retraining on a group-aware split with a
counterfeit-family hold-out is the most important next step.

%% ======================================================================
\section{Conclusion}\label{sec:conclusion}

Our smartphone system combines multi-view fusion, a three-way decision, an OOD gate, and
Grad-CAM, runs in real time on a CPU, and reaches 98.08\% note-level accuracy (AUC
0.994) on JaalTaka. The audit shows that label conflicts among duplicated notes explain
every error and that on duplicate-free notes even simple baselines are error-free.
Progress now depends on evaluation data that separates counterfeit families, sessions,
and devices.

%% ======================================================================
\backmatter

\bmhead{Acknowledgements}
The authors thank the Department of Software Engineering, Daffodil International
University, for its support and the JaalTaka creators for releasing their data.

\section*{Declarations}

\begin{itemize}
\item \textbf{Funding:} This research received no external funding.
\item \textbf{Competing interests:} The authors declare no competing interests.
\item \textbf{Ethics approval and consent to participate:} Not applicable. This study
involved no human participants, personal data, or animals; it used only banknote
photographs from a public CC~BY~4.0 dataset~\cite{jaaltaka_data}.
\item \textbf{Consent for publication:} Not applicable.
\item \textbf{Data availability:} JaalTaka is available from Mendeley
Data~\cite{jaaltaka_data}. The duplicate-group list is available from the corresponding
author and will be released publicly upon publication.
\item \textbf{Materials availability:} Not applicable.
\item \textbf{Code availability:} The model and code are available from the
corresponding author and will be released publicly upon publication.
\item \textbf{Author contributions:} D.~Biswas: conceptualisation, software,
experiments, analysis, and original draft. I.~Mahmud: supervision, methodology review,
and manuscript editing.
%% TODO (authors): adapt this statement to the journal's policy and to how AI tools
%% were actually used.
\item \textbf{Use of generative AI:} The authors used an AI assistant to help write
analysis scripts and edit the text, and checked all code, results, and references.
\end{itemize}

%% ----------------------------------------------------------------------
\bibliography{references}

\end{document}
